\section{Implementation}


\subsection{Go Idioms for Cloud Services}

Our implementation leverages Go idioms suited to cloud environments:

\textbf{Context propagation.} All functions accept \texttt{context.Context} as the first parameter, enabling deadline propagation and cancellation.

\textbf{Interface-based design.} Services depend on interfaces, not concrete implementations, facilitating testing and substitution.

\textbf{Error wrapping.} Errors include context using \texttt{fmt.Errorf} with the \texttt{\%w} verb, enabling error chain inspection.

\subsection{Testing Strategy}

We employ three testing levels:

\begin{enumerate}
\item \textbf{Unit tests}: Pure functions with mocked dependencies.
\item \textbf{Integration tests}: Full service tests against local Datastore emulator.
\item \textbf{End-to-end tests}: Complete flows against staging environment.
\end{enumerate}

The local development server replicates production behavior:

\begin{lstlisting}[language=bash,caption=Local development]
$ dev_appserver.py app.yaml
INFO: Starting module "default" at http://localhost:8080
INFO: Starting admin server at http://localhost:8000
\end{lstlisting}

\subsection{Deployment Pipeline}

Deployments follow a progressive rollout strategy:

\begin{enumerate}
\item Push to \texttt{main} triggers CI/CD pipeline.
\item Tests execute against emulated services.
\item Passing builds deploy to staging (1\% traffic).
\item Automated canary analysis compares error rates.
\item Gradual traffic shift to 100\% over 30 minutes.
\end{enumerate}

